\clearpage\thispagestyle{empty}
\pdfbookmark[0]{Quellen zu Gauß an Bode (5. August 1825)}{bodeaug1825-sources}
\begin{center}\textbf{Quellen zu Gauß an Bode (5. August 1825)}\end{center}
\noindent Diese KI-Übertragung ins heutige Deutsch folgt Gauß’ Brief vom 5. August 1825 in \href{https://gdz.sub.uni-goettingen.de/id/PPN236007467}{den \textit{Werken}, Band~VI} (Göttingen, 1874), S.~451. Johann Elert Bodes \textit{Astronomisches Jahrbuch für 1828} (Berlin, 1825), \href{https://books.google.nl/books?id=GDAPAAAAQAAJ&pg=PA171&redir_esc=y}{S.~171} und \href{https://books.google.nl/books?id=GDAPAAAAQAAJ&pg=PA172&redir_esc=y}{S.~172}, druckt den Beitrag unter der Überschrift \textit{Meridian-Beobachtungen der Pallas und Ceres im Jahr 1825 vom Herrn Hofrath Ritter Gauss in Göttingen}. Die Links führen zum Digitalisat in Google Books.
\par\smallskip\noindent Das Jahrbuch nennt den 5. August als Einsendedatum. In der Werkausgabe steht das Datum in der Orts- und Datumszeile des Briefs aus Göttingen. Beide Drucke enthalten vierzehn Beobachtungen der Pallas und sieben Beobachtungen der Ceres. KI-Übertragung und KI-Satz.
\clearpage
